\documentclass[letterpaper,twocolumn,10pt]{article}
\usepackage{usenix2019_v3}

\usepackage{tikz}
\usepackage{amsmath}
\usepackage{filecontents}

\begin{document}

\title{\Large \bf Karna: A High-Performance, Tiered Storage System for Modern Workloads}

\author{
{\rm Author Name}\\
Institution
} % end author

\maketitle

\begin{abstract}
Modern data-intensive workloads demand high throughput, low latency, and massive scalability. We introduce Karna, a novel tiered storage architecture that intelligently migrates hot and cold data across storage layers using predictive machine learning models, optimizing for both performance and cost.
\end{abstract}

\section{Introduction}
With the advent of NVMe, persistent memory, and scalable cloud storage, managing data placement across diverse storage tiers is a critical challenge. Existing heuristic-based approaches fail to adapt to dynamic workload patterns. Karna introduces an ML-driven data placement strategy.

\section{Methodology}
\subsection{Tiered Storage Architecture}
Karna abstracts underlying storage media into a unified address space. Data is segmented into variable-sized chunks, and placement decisions are made asynchronously.
\subsection{Predictive Migration}
We deploy a lightweight recurrent neural network to predict access frequencies, allowing proactive data staging to faster tiers before spikes in demand.

\section{Experiments}
\subsection{Setup}
We benchmark Karna against state-of-the-art tiered file systems using standard I/O traces (e.g., FIU, MSR).
\subsection{Results}
% Placeholder results
Karna achieves up to 40\% reduction in average read latency and a 25\% improvement in overall throughput under highly skewed workload conditions.

\section{Conclusion}
Karna demonstrates that integrating machine learning into the storage data path can yield significant performance benefits for tiered storage systems.

\bibliographystyle{plain}
% \bibliography{refs}

\end{document}
